% ======================================================================================
% BOOK MACROS
% ======================================================================================
% This file is the book's own macro file: the PDF gets it by \input, and
% build/macros.py parses the same file into the MathJax macro table the website
% uses, so the two formats always agree. Add a macro here (\def, \newcommand,
% \providecommand, \DeclareMathOperator) and both editions have it.
%
% WHAT IS IN HERE: only what any mathematics book wants -- delimiters, a handful
% of common operators, the blackboard and calligraphic alphabets, bold Latin
% letters for vectors and matrices, derivatives and limits, and the two-slot
% bracket pattern \coord / \mtx.
%
% WHAT IS NOT: the subject-specific macros the linear algebra book accumulated
% (its operators, bold Greek vectors, transpose and adjoint shorthands, Dirac
% notation, the probability and distribution sets). Those sit unloaded in
% latex/macros-linear-algebra.tex.example. Raid that file for what this book
% needs and paste it in here; do not \input it wholesale, since a macro nobody
% uses is a name this book cannot have.
%
% Two things to know before editing:
%   - Several of these names belong to LaTeX in text mode (\b \c \d \i \j \l \o
%     \t \u \v \H \L \O \P \S are accents or symbols). Redefining them is
%     deliberate and long-standing; latex/preamble.tex restores \H after
%     hyperref, which reclaims it. If a new macro shadows a text command the
%     book actually uses, that is on you.
%   - \providecommand here is a no-op in LaTeX when latex/latex-template.sty has
%     already defined the name, but it is still what the website reads. Keep the
%     two bodies saying the same thing.
% ======================================================================================

% ======================================================================================
% NUMBER SETS
% ======================================================================================
% \mathbb (amssymb/msbm10), not bbm's \mathbb: bbm ships no Type 1 face in common TeX
% installations, so pdflatex falls back to a bitmap font. That put Type 3 glyphs in every
% PDF (blurry, unsearchable, uncopyable) and turned blackboard letters inside TikZ pictures
% into masked rectangles in the SVGs.
% \ensuremath so the macros work in text mode too (e.g. in an environment title).
\def \nA {\ensuremath{\mathbb{A}}}
\def \nB {\ensuremath{\mathbb{B}}}
\def \nC {\ensuremath{\mathbb{C}}}
\def \nD {\ensuremath{\mathbb{D}}}
\def \nE {\ensuremath{\mathbb{E}}}
\def \nF {\ensuremath{\mathbb{F}}}
\def \nG {\ensuremath{\mathbb{G}}}
\def \nH {\ensuremath{\mathbb{H}}}
\def \nI {\ensuremath{\mathbb{I}}}
\def \nJ {\ensuremath{\mathbb{J}}}
\def \nK {\ensuremath{\mathbb{K}}}
\def \nL {\ensuremath{\mathbb{L}}}
\def \nM {\ensuremath{\mathbb{M}}}
\def \nN {\ensuremath{\mathbb{N}}}
\def \nO {\ensuremath{\mathbb{O}}}
\def \nP {\ensuremath{\mathbb{P}}}
\def \nQ {\ensuremath{\mathbb{Q}}}
\def \nR {\ensuremath{\mathbb{R}}}
\def \nS {\ensuremath{\mathbb{S}}}
\def \nT {\ensuremath{\mathbb{T}}}
\def \nU {\ensuremath{\mathbb{U}}}
\def \nV {\ensuremath{\mathbb{V}}}
\def \nW {\ensuremath{\mathbb{W}}}
\def \nX {\ensuremath{\mathbb{X}}}
\def \nY {\ensuremath{\mathbb{Y}}}
\def \nZ {\ensuremath{\mathbb{Z}}}

% ======================================================================================
% CALLIGRAPHIC LETTERS
% ======================================================================================
% Families, collections, ordered bases, sigma-algebras: \sB, \sF, ...
\def \sA {\mathcal{A}}
\def \sB {\mathcal{B}}
\def \sC {\mathcal{C}}
\def \sD {\mathcal{D}}
\def \sE {\mathcal{E}}
\def \sF {\mathcal{F}}
\def \sG {\mathcal{G}}
\def \sH {\mathcal{H}}
\def \sI {\mathcal{I}}
\def \sJ {\mathcal{J}}
\def \sK {\mathcal{K}}
\def \sL {\mathcal{L}}
\def \sM {\mathcal{M}}
\def \sN {\mathcal{N}}
\def \sO {\mathcal{O}}
\def \sP {\mathcal{P}}
\def \sQ {\mathcal{Q}}
\def \sR {\mathcal{R}}
\def \sS {\mathcal{S}}
\def \sT {\mathcal{T}}
\def \sU {\mathcal{U}}
\def \sV {\mathcal{V}}
\def \sW {\mathcal{W}}
\def \sX {\mathcal{X}}
\def \sY {\mathcal{Y}}
\def \sZ {\mathcal{Z}}

% ======================================================================================
% BOLD LATIN LETTERS
% ======================================================================================
% The book-wide convention both existing books follow: bold lowercase for
% vectors, bold uppercase for matrices, plain italic for maps and spaces.
\def \a {\mathbf{a}}
\def \b {\mathbf{b}}
\def \c {\mathbf{c}}
\def \d {\mathbf{d}}
\def \e {\mathbf{e}}
\def \f {\mathbf{f}}
\def \g {\mathbf{g}}
\def \h {\mathbf{h}}
\def \i {\mathbf{i}}
\def \j {\mathbf{j}}
\def \k {\mathbf{k}}
\def \l {\mathbf{l}}
\def \m {\mathbf{m}}
\def \n {\mathbf{n}}
\def \o {\mathbf{o}}
\def \p {\mathbf{p}}
\def \q {\mathbf{q}}
\def \r {\mathbf{r}}
\def \s {\mathbf{s}}
\def \t {\mathbf{t}}
\def \u {\mathbf{u}}
\def \v {\mathbf{v}}
\def \w {\mathbf{w}}
\def \x {\mathbf{x}}
\def \y {\mathbf{y}}
\def \z {\mathbf{z}}
\def \A {\mathbf{A}}
\def \B {\mathbf{B}}
\def \C {\mathbf{C}}
\def \D {\mathbf{D}}
\def \E {\mathbf{E}}
\def \F {\mathbf{F}}
\def \G {\mathbf{G}}
\def \H {\mathbf{H}}
\def \I {\mathbf{I}}
\def \J {\mathbf{J}}
\def \K {\mathbf{K}}
\def \L {\mathbf{L}}
\def \M {\mathbf{M}}
\def \N {\mathbf{N}}
\def \O {\mathbf{O}}
\def \P {\mathbf{P}}
\def \Q {\mathbf{Q}}
\def \R {\mathbf{R}}
\def \S {\mathbf{S}}
\def \T {\mathbf{T}}
\def \U {\mathbf{U}}
\def \V {\mathbf{V}}
\def \W {\mathbf{W}}
\def \X {\mathbf{X}}
\def \Y {\mathbf{Y}}
\def \Z {\mathbf{Z}}

\def\0{\mathbf{0}}
\def\1{\mathbf{1}}

% ======================================================================================
% DELIMITERS
% ======================================================================================
\providecommand{\abs}[1]{\left| #1 \right|}
\providecommand{\norm}[1]{\left\| #1 \right\|}
\providecommand{\card}[1]{\left| #1 \right|}
\providecommand{\powerset}[1]{\mathcal{P}\left( #1 \right)}
\providecommand{\inner}[2]{\langle #1, #2 \rangle}
\providecommand{\ip}[2]{\ensuremath{\langle#1,#2\rangle}}
\providecommand{\conj}[1]{\overline{#1}}
\providecommand{\closure}[1]{\overline{#1}}

% ======================================================================================
% THE TWO-SLOT BRACKET PATTERN
% ======================================================================================
% \coord{x}{B} renders [x]_B and \mtx{T}{B}{C} renders [T]_B^C: an object in
% brackets, with the data it is taken relative to hanging off the brackets. The
% linear algebra book uses it for coordinate vectors and for the matrix of a map
% relative to a source and a target basis, and NOTATION.md fixes the slot order
% once so no section invents its own. Any book with "X relative to Y" notation
% wants the same discipline; rename the macros if the words change.
\providecommand{\coord}[2]{[#1]_{#2}}
\providecommand{\mtx}[3]{[#1]_{#2}^{#3}}

% ======================================================================================
% COMMON OPERATORS
% ======================================================================================
\def\id{\operatorname{id}}
\def\sgn{\operatorname{sgn}}

% ======================================================================================
% DERIVATIVES, LIMITS, DIFFERENTIALS
% ======================================================================================
\def \dd {\mathrm{d}}
\providecommand*{\diff}{\mathop{}\:\!\mathrm d}
\providecommand{\pd}[2]{\frac{\partial #1}{\partial #2}}
\providecommand{\pdd}[2]{\frac{\partial^2 #1}{\partial #2^2}}
\providecommand{\td}[2]{\frac{d #1}{d #2}}
\providecommand{\tdd}[2]{\frac{d^2 #1}{d #2^2}}
\providecommand{\limi}[1]{\lim_{#1 \to \infty}}
\providecommand{\limz}[1]{\lim_{#1 \to 0}}

% ======================================================================================
% MISCELLANY
% ======================================================================================
% "is defined to be". \coloneqq itself comes from mathtools in the PDF and from
% build/macros.py's WEB_DEFAULTS on the web.
\providecommand{\defeq}{\stackrel{\smash{\textnormal{\tiny def}}}{=}}
% Topology
\providecommand{\interior}[1]{\text{int}\left( #1 \right)}
\providecommand{\boundary}[1]{\partial #1}
% One-letter alphabet switches, for the odd letter that has no macro of its own
\providecommand{\bb}[1]{\mathbb{#1}}
\providecommand{\cl}[1]{\mathcal{#1}}
\providecommand{\rom}[1]{\mathrm{#1}}
% A fallback only: the PDF has the real \bm from the bm package, and
% build/macros.py's WEB_OVERRIDES gives the website \boldsymbol to match it.
\providecommand{\bm}[1]{\mathbf{#1}}
